% Lessons 63--75: DP, validation, tries, search, and the final Blind 75 topics.

\patternintoc{Grid DP}
\problemstart{Unique Paths}{62}{Medium}{Grid dynamic programming}{One-dimensional DP}

\subsection{The Problem}
Count paths from the top-left to bottom-right of an \(m\times n\) grid when
each move goes only right or down.

\subsection{Each Cell Receives Paths from Two Directions}
The number of ways into a cell is the sum of ways into the cell above and the
cell to the left. A one-dimensional row stores both: before updating
\texttt{dp[c]} is the value from above, while \texttt{dp[c-1]} is from left.

\begin{invariantbox}{dp[c] is the path count for the current processed row}
After column \(c\) is updated, \texttt{dp[c]} contains all paths to the current
cell. Entries to its right still represent the previous row.
\end{invariantbox}

\begin{visualbox}{Paths flow from above and left}
\centering
\begin{tikzpicture}[font=\small,>={Stealth[length=2.2mm,width=1.55mm]}]
  \matrix[matrix of nodes,nodes={arraycell},row sep=-\pgflinewidth,column sep=-\pgflinewidth] (m) {1&1&1\\1&2&3\\1&3&6\\};
  \node[draw=teal,very thick,fit=(m-2-3),inner sep=1mm] {};
  \node[draw=orange,very thick,fit=(m-3-2),inner sep=1mm] {};
  \node[draw=navy,very thick,fit=(m-3-3),inner sep=1mm] {};
  \node[right=10mm of m-2-3,align=left]
    {from above: 3\\from left: 3\\new value: $3+3=6$};
\end{tikzpicture}
\end{visualbox}

\subsection{Pseudocode}
\begin{lstlisting}[style=pseudocode]
dp[0...n-1] = 1
repeat for each remaining row:
    for column from 1 to n-1:
        dp[column] += dp[column - 1]
return dp[n - 1]
\end{lstlisting}

\subsection{C++ Implementation}
\begin{lstlisting}[style=compactcode,float=htbp,caption={Complete solution: Unique Paths}]
#include <vector>

using namespace std;

class Solution {
public:
    int uniquePaths(int m, int n) {
        vector<int> dp(n, 1);
        for (int r = 1; r < m; ++r)
            for (int c = 1; c < n; ++c) dp[c] += dp[c - 1];
        return dp[n - 1];
    }
};
\end{lstlisting}

\subsection{Time and Space}
Time is \(O(mn)\), space is \(O(n)\). The first row and first column contain
one path each, which explains the initialization to one.

\lessonexpansion
  {For a \(3\times3\) grid, start \([1,1,1]\). The second row updates to
   \([1,2,3]\), and the third row becomes \([1,3,6]\). The final entry shows
   six ways to reach the lower-right corner.}
  {Every cell can be entered only from above or from the left. During a
   left-to-right update, \texttt{dp[c]} still stores the count from above and
   \texttt{dp[c-1]} stores the updated count from the left, so their sum is
   exactly the new state.}
  {Check \(1\times1\), one row, one column, a square grid, and highly
   rectangular grids.}
  {I compress the two-dimensional update rule into one row. Each entry adds
   the ways from above and left, giving \(O(mn)\) time and \(O(n)\) space.}

\chaptercheckpoint{Count monotone paths through a rectangular grid.}
  {One DP row updated from above and left.}
  {Initializing the starting boundary to zero.}{$O(mn)$ time and $O(n)$ space.}


\patternintoc{String Scanning and Stacks}
\problemstart{Valid Anagram}{242}{Easy}{Character-frequency equality}{Count array}

\subsection{The Problem}
Return whether two strings contain exactly the same characters with the same
multiplicities.

\begin{invariantbox}{the count array stores the remaining character imbalance}
Incrementing for the first string and decrementing for the second leaves all
zeros exactly when their multisets are identical.
\end{invariantbox}

\begin{visualbox}{Matching letters cancel in the frequency table}
\centering
\begin{tikzpicture}[font=\small]
  \node[dsnode] (a) {\texttt{anagram}\\add counts};
  \node[dsnode,below=9mm of a] (b) {\texttt{nagaram}\\subtract counts};
  \node[dsnode,active,right=25mm of a,yshift=-5mm] (z) {all 26\\entries zero};
  \draw[flowarrow] (a)--(z); \draw[flowarrow,teal] (b)--(z);
\end{tikzpicture}
\end{visualbox}

\subsection{Pseudocode}
\begin{lstlisting}[style=pseudocode]
if lengths differ: return false
counts = 26 zeroes
for i from 0 to n-1:
    increment counts[s[i]]
    decrement counts[t[i]]
return every count is zero
\end{lstlisting}

\subsection{C++ Implementation}
\begin{lstlisting}[style=compactcode,float=htbp,caption={Complete solution: Valid Anagram}]
#include <array>
#include <string>

using namespace std;

class Solution {
public:
    bool isAnagram(string s, string t) {
        if (s.size() != t.size()) return false;
        array<int,26> count{};
        for (int i=0;i<(int)s.size();++i){++count[s[i]-'a'];--count[t[i]-'a'];}
        for(int value:count) if(value!=0) return false;
        return true;
    }
};
\end{lstlisting}

\subsection{Time and Space}
Time is \(O(n)\), alphabet space is \(O(1)\). Length mismatch can be rejected
immediately.

\lessonexpansion
  {For \texttt{anagram} and \texttt{nagaram}, incrementing counts from the
   first word and decrementing them with the second returns every letter count
   to zero. For \texttt{rat} and \texttt{car}, the \texttt{t} and
   \texttt{c} entries remain nonzero.}
  {Equal-length lowercase strings are anagrams exactly when every character
   occurs the same number of times. The combined increment/decrement array
   measures those frequency differences directly.}
  {Check empty strings, one character, repeated letters, a length mismatch,
   and strings differing by only one occurrence.}
  {I compare character multisets with one fixed-size count array. Matching
   additions and removals cancel; all zeroes at the end prove an anagram.}

\chaptercheckpoint{Compare two lowercase-character multisets.}
  {A fixed frequency-difference array.}
  {Using a set and losing duplicate counts.}{$O(n)$ time and $O(1)$ space.}


\problemstart{Valid Palindrome}{125}{Easy}{Filtered two pointers}{String}

\subsection{The Problem}
Ignoring non-alphanumeric characters and letter case, determine whether a
string is a palindrome.

\begin{invariantbox}{all characters outside the pointers already match}
Each pointer skips irrelevant characters. When both point to alphanumeric
characters, their lowercase forms must match before the search moves inward.
\end{invariantbox}

\begin{visualbox}{Pointers ignore punctuation and compare useful characters}
\centering
\begin{tikzpicture}[font=\small,>={Stealth[length=2.2mm,width=1.55mm]}]
  \node[draw=navy,fill=softblue,rounded corners,inner sep=3mm] (s)
    {\texttt{A man, a plan, a canal: Panama}};
  \node[dsnode,fill=softteal,below=7mm of s] (clean)
    {keep lowercase\\alphanumeric characters};
  \draw[flowarrow,teal] (s)--(clean);

  \coordinate (rowcenter) at ([yshift=-12mm]clean.south);
  \foreach \ch [count=\i from 0] in
    {a,m,a,n,a,p,l,a,n,a,c,a,n,a,l,p,a,n,a,m,a} {
      \node[arraycell,minimum width=5mm,minimum height=7mm,
        inner sep=1pt,font=\scriptsize] (c\i)
        at ([xshift={(-50+5*\i)*1mm}]rowcenter) {\ch};
  }
  \node[draw=teal,line width=1pt,fit=(c0),inner sep=0.5pt] {};
  \node[draw=orange,line width=1pt,fit=(c20),inner sep=0.5pt] {};

  \node[below=9mm of c0,text=teal,font=\bfseries] (left) {left};
  \node[below=9mm of c20,text=orange,font=\bfseries] (right) {right};
  \draw[flowarrow,teal] (left.north)--(c0.south);
  \draw[flowarrow,orange] (right.north)--(c20.south);
  \node[below=25mm of rowcenter,text=navy]
    {the end characters match: \texttt{a} = \texttt{a}; move both pointers inward};
\end{tikzpicture}
\end{visualbox}

\subsection{Pseudocode}
\begin{lstlisting}[style=pseudocode]
left = 0; right = n - 1
while left < right:
    skip non-alphanumeric characters on both sides
    if lowercase characters differ: return false
    move both boundaries inward
return true
\end{lstlisting}

\subsection{C++ Implementation}
\begin{lstlisting}[style=compactcode,float=htbp,caption={Complete solution: Valid Palindrome}]
#include <cctype>
#include <string>

using namespace std;

class Solution {
public:
    bool isPalindrome(string s) {
        int left=0,right=(int)s.size()-1;
        while(left<right){
            while(left<right && !isalnum((unsigned char)s[left])) ++left;
            while(left<right && !isalnum((unsigned char)s[right])) --right;
            if(tolower((unsigned char)s[left])!=tolower((unsigned char)s[right])) return false;
            ++left;--right;
        }
        return true;
    }
};
\end{lstlisting}

\subsection{Time and Space}
Each pointer moves inward once: \(O(n)\) time and \(O(1)\) space. Casting to
unsigned char makes the classification functions safe.

\lessonexpansion
  {For \texttt{"A man, a plan, a canal: Panama"}, the pointers skip spaces
   and punctuation, compare \texttt{a} with \texttt{a}, then continue inward
   case-insensitively until they cross.}
  {Skipping non-alphanumeric characters makes the two pointers identify the
   next pair in the filtered string. If every mirrored pair matches after case
   normalization, that filtered string equals its reverse.}
  {Check empty text, punctuation only, mixed case, digits, an early mismatch,
   and a mismatch near the center.}
  {I scan inward from both ends, ignoring characters outside the problem's
   alphabet and comparing the remaining characters without case.}

\chaptercheckpoint{Compare text after ignoring punctuation and case.}
  {Two inward pointers.}
  {Passing negative signed chars directly to cctype functions.}{$O(n)$ time.}


\problemstart{Valid Parentheses}{20}{Easy}{Delimiter matching}{Stack}

\subsection{The Problem}
Return whether every opening bracket is closed by the correct type in the
correct order.

\begin{invariantbox}{the stack contains exactly the unmatched openings}
An opening bracket is pushed. A closing bracket must match the stack top,
because that is the most recent still-unclosed delimiter.
\end{invariantbox}

\begin{visualbox}{Nested delimiters close in LIFO order}
\centering
\begin{tikzpicture}[font=\small]
  \node[draw=navy,fill=softblue,minimum width=25mm] (a) {(};
  \node[draw=navy,fill=softblue,minimum width=25mm,above=0mm of a] (b) {[};
  \node[draw=orange,fill=softorange,minimum width=25mm,above=0mm of b] (c) {\{};
  \node[right=12mm of c,text=orange]{next \texttt{\}} pops top};
\end{tikzpicture}
\end{visualbox}

\subsection{Pseudocode}
\begin{lstlisting}[style=pseudocode]
stack = empty
for character in string:
    if opening delimiter: push it
    else if stack empty or top is not its partner: return false
    else pop stack
return stack is empty
\end{lstlisting}

\subsection{C++ Implementation}
\begin{lstlisting}[style=compactcode,float=htbp,caption={Complete solution: Valid Parentheses}]
#include <stack>
#include <string>

using namespace std;

class Solution {
public:
    bool isValid(string s) {
        stack<char> openings;
        for(char ch:s){
            if(ch=='('||ch=='['||ch=='{') openings.push(ch);
            else {
                if(openings.empty()) return false;
                char open=openings.top(); openings.pop();
                if((ch==')'&&open!='(')||(ch==']'&&open!='[')||(ch=='}'&&open!='{')) return false;
            }
        }
        return openings.empty();
    }
};
\end{lstlisting}

\subsection{Time and Space}
Time is \(O(n)\), stack space \(O(n)\). An empty stack before a closing bracket
and leftover openings after the scan are both invalid.

\lessonexpansion
  {For \texttt{([\{\}])}, push each opening delimiter. The closing brace pops
   \texttt{\{}, the closing bracket pops \texttt{[}, and the closing
   parenthesis pops \texttt{(}. For \texttt{([)]}, the first closing
   parenthesis conflicts with the stack top \texttt{[}.}
  {Nested delimiters must close in reverse opening order, exactly the order
   supplied by a stack. Every closing delimiter is validated against the only
   opening delimiter it could legally match.}
  {Check an empty string, one unmatched delimiter, adjacent pairs, deep
   nesting, crossed types, and leftover opening delimiters.}
  {I push opening delimiters and require every closing delimiter to match the
   stack top. The stack must also be empty after the complete scan.}

\chaptercheckpoint{Nested pairs must close in reverse opening order.}
  {A stack of unmatched opening brackets.}
  {Checking only total bracket counts.}{$O(n)$ time and $O(n)$ space.}


\patternintoc{Trees and Word Segmentation}
\problemstart{Validate Binary Search Tree}{98}{Medium}{Recursive value bounds}{Binary tree}

\subsection{The Problem}
Return whether every node satisfies the global binary-search-tree ordering.

\subsection{Carry Ancestor Bounds}
Local child comparisons are inenough. Each recursive call receives the
strict lower and upper bounds imposed by all ancestors.

\begin{invariantbox}{every node in the current subtree must lie inside its bounds}
The left call tightens the upper bound to the current value; the right call
tightens the lower bound. Strict inequalities reject duplicates.
\end{invariantbox}

\begin{visualbox}{Ancestor bounds reach deeper than the parent}
\centering
\begin{tikzpicture}[font=\small]
  \node[trienode] (r) {5}; \node[trienode,below left=11mm and 15mm of r] (a) {1};
  \node[trienode,below right=11mm and 15mm of r] (b) {7};
  \node[trienode,active,below left=10mm and 7mm of b] (c) {4};
  \draw[-,navy] (r)--(a) (r)--(b) (b)--(c);
  \node[right=8mm of c,text=orange]{invalid: right subtree must be $>5$};
\end{tikzpicture}
\end{visualbox}

\subsection{Pseudocode}
\begin{lstlisting}[style=pseudocode]
valid(node, lower, upper):
    if node is null: return true
    if node.value <= lower or node.value >= upper: return false
    return valid(node.left, lower, node.value)
           AND valid(node.right, node.value, upper)
\end{lstlisting}

\subsection{C++ Implementation}
\begin{lstlisting}[style=compactcode,float=htbp,caption={Complete solution: Validate BST}]
#include <limits>

using namespace std;

class Solution {
    bool valid(TreeNode* node,long long low,long long high){
        if(node==nullptr) return true;
        if(node->val<=low||node->val>=high) return false;
        return valid(node->left,low,node->val)&&valid(node->right,node->val,high);
    }
public:
    bool isValidBST(TreeNode* root){return valid(root,numeric_limits<long long>::lowest(),numeric_limits<long long>::max());}
};
\end{lstlisting}

\subsection{Time and Space}
Every node is checked once: \(O(n)\) time, \(O(h)\) stack space. Wide bounds
avoid collisions with valid \texttt{int} endpoints.

\lessonexpansion
  {In \texttt{[5,1,4,null,null,3,6]}, node 3 is locally smaller than parent 4
   but lies in the root's right subtree, where every value must exceed 5.
   Passing the inherited lower bound 5 exposes the violation.}
  {Each recursive call carries the complete range allowed by all ancestors.
   Restricting the upper bound on a left edge and the lower bound on a right
   edge makes every accepted node satisfy the global BST definition.}
  {Check valid integer endpoints, duplicate values, a deep ancestor violation,
   a skewed valid BST, and an empty tree.}
  {I validate each node against inherited lower and upper bounds rather than
   only its parent. Child calls tighten exactly one side of that range.}

\chaptercheckpoint{BST the valid condition depends on all ancestors, not only parents.}
  {Strict lower and upper bounds.}
  {Accepting a deep value that violates an ancestor.}{$O(n)$ time and $O(h)$ space.}


\problemstart{Word Break}{139}{Medium}{Prefix segmentation DP}{String and hash set}

\subsection{The Problem}
Return whether a string can be split into one or more dictionary words. Words
may be reused.

\begin{invariantbox}{dp[i] means the prefix ending before i is segmentable}
For every true \texttt{dp[i]}, trying dictionary words tests whether one legal
word can extend that already valid prefix. So every true state corresponds
to a complete segmentation.
\end{invariantbox}

\begin{visualbox}{True prefix states jump across dictionary words}
\centering
\begin{tikzpicture}[font=\small,>={Stealth[length=2.2mm,width=1.55mm]}]
  \foreach \x/\v in {0/T,1/F,2/F,3/F,4/T,5/F,6/F,7/F,8/T}{\node[arraycell] (d\x) at (0.8*\x,0) {\v};}
  \draw[->,very thick,teal] (d0.north) to[bend left] node[edgelabel,above]{\texttt{leet}} (d4.north);
  \draw[->,very thick,orange] (d4.north) to[bend left] node[edgelabel,above]{\texttt{code}} (d8.north);
\end{tikzpicture}
\end{visualbox}

\subsection{Pseudocode}
\begin{lstlisting}[style=pseudocode]
dp[0] = true
for end from 1 to n:
    for start from 0 to end-1:
        if dp[start] and s[start...end) is in dictionary:
            dp[end] = true; break
return dp[n]
\end{lstlisting}

\subsection{C++ Implementation}
\begin{lstlisting}[style=compactcode,float=htbp,caption={Complete solution: Word Break}]
#include <string>
#include <unordered_set>
#include <vector>

using namespace std;

class Solution {
public:
    bool wordBreak(string s,vector<string>& words){
        unordered_set<string> dict(words.begin(),words.end());
        vector<bool> dp(s.size()+1,false); dp[0]=true;
        for(int end=1;end<=(int)s.size();++end)
            for(int start=0;start<end;++start)
                if(dp[start]&&dict.count(s.substr(start,end-start))){dp[end]=true;break;}
        return dp[s.size()];
    }
};
\end{lstlisting}

\subsection{Time and Space}
The direct prefix DP uses \(O(n^2)\) transitions plus substring construction
cost; space is \(O(n+D)\). Limiting starts by maximum word length improves the
constant factors.

\lessonexpansion
  {For \texttt{leetcode} with dictionary \texttt{[leet,code]}, begin with
   \(dp[0]=true\). The word \texttt{leet} makes \(dp[4]\) true. From that
   reachable boundary, \texttt{code} makes \(dp[8]\) true, proving the full
   string can be segmented.}
  {A true \(dp[end]\) is created only by a true earlier boundary followed by a
   dictionary word, so every true state has a valid segmentation. The reverse is also true:,
   the final word of any valid segmentation begins at some true earlier
   boundary, so the transition discovers every possible solution.}
  {Check an empty prefix, one dictionary word, repeated word use, overlapping
   choices such as \texttt{cars}, a long almost-valid suffix, and an input
   with no valid first word.}
  {I use prefix dynamic programming because a locally valid word choice may
   block the suffix. Each true index records that the entire prefix is valid,
   allowing all later dictionary endings to build on it.}

\chaptercheckpoint{A string must be partitioned into reusable dictionary words.}
  {Boolean states for segmentable prefixes.}
  {Greedily taking the first matching word.}{$O(n^2)$ DP transitions.}


\patternintoc{Trie-Guided Grid Search}
\problemstart{Word Search II}{212}{Hard}{Trie-pruned backtracking}{Trie and grid}

\subsection{The Problem}
Find all dictionary words that can be formed by adjacent grid cells without
reusing a cell within one word.

\subsection{Share Prefix Work with a Trie}
Insert all words into a Trie. DFS from each cell follows only Trie edges that
exist; impossible prefixes stop immediately. Store a complete word at terminal
nodes and clear it after reporting to avoid duplicates.

\begin{invariantbox}{the DFS path equals the Trie prefix at its node}
Every active board path uses distinct marked cells and spells exactly the
characters from the Trie root to the current Trie node. Missing edges safely
prune every dictionary word sharing that impossible prefix.
\end{invariantbox}

\begin{visualbox}{Board paths and Trie prefixes advance together}
\centering
\begin{tikzpicture}[font=\small]
  \matrix[matrix of nodes,nodes={arraycell},row sep=-\pgflinewidth,
    column sep=-\pgflinewidth] (b) {
    |[fill=softorange,draw=orange,line width=1pt]| o &
    |[fill=softorange,draw=orange,line width=1pt]| a & a & n\\
    e & |[fill=softorange,draw=orange,line width=1pt]| t & a & e\\
    i & |[fill=softorange,draw=orange,line width=1pt]| h & k & r\\
    i & f & l & v\\};
  \node[right=20mm of b,dsnode] {board path: \texttt{oath}\\
    same Trie prefix: \(o-a-t-h\)};
\end{tikzpicture}
\end{visualbox}

\subsection{Pseudocode}
\begin{lstlisting}[style=pseudocode]
insert every word into a trie
for every board cell:
    DFS(cell, trie root)
DFS(cell, trie node):
    stop if outside, visited, or character edge is absent
    move to child; report and clear any completed word
    mark cell, explore four neighbors, then restore cell
return reported words
\end{lstlisting}

\subsection{C++ Implementation}
\begin{lstlisting}[style=compactcode,float=htbp,caption={Complete solution: Word Search II}]
#include <array>
#include <string>
#include <vector>

using namespace std;

class Solution {
    struct Trie { array<Trie*,26> next{}; string word; };
    void search(vector<vector<char>>& b,int r,int c,Trie* node,vector<string>& out){
        if(r<0||c<0||r==(int)b.size()||c==(int)b[0].size()||b[r][c]=='#') return;
        char ch=b[r][c]; Trie* child=node->next[ch-'a']; if(child==nullptr) return;
        if(!child->word.empty()){out.push_back(child->word);child->word.clear();}
        b[r][c]='#'; search(b,r+1,c,child,out);search(b,r-1,c,child,out);
        search(b,r,c+1,child,out);search(b,r,c-1,child,out); b[r][c]=ch;
    }
public:
    vector<string> findWords(vector<vector<char>>& board,vector<string>& words){
        Trie* root=new Trie();
        for(const auto& word:words){Trie* p=root;for(char ch:word){if(!p->next[ch-'a'])p->next[ch-'a']=new Trie();p=p->next[ch-'a'];}p->word=word;}
        vector<string> out; for(int r=0;r<(int)board.size();++r)for(int c=0;c<(int)board[0].size();++c)search(board,r,c,root,out); return out;
    }
};
\end{lstlisting}

\subsection{Time, Space, and Safety}
Trie construction is \(O(W)\) for total dictionary characters. Search is
bounded by board starts and valid prefix branches; the loose worst case is
\(O(mn\,4^L)\). Backtracking restores each cell. Production code should own
Trie nodes with smart pointers or a node pool.

\lessonexpansion
  {On the standard board, DFS beginning at \texttt{o} follows Trie edges
   \texttt{o->a->t->h} and reports \texttt{oath}. A neighboring character with
   no Trie edge stops immediately. Clearing the terminal word prevents another
   board path from returning the same dictionary word twice.}
  {The active board path and Trie path always spell the same prefix. So
   a missing Trie child proves that no dictionary word can extend that board
   path. Every reported terminal node corresponds to a valid non-reusing path,
   and every valid dictionary path is explored from its first cell.}
  {Check duplicate dictionary entries, words sharing long prefixes, a one-cell
   board, a word longer than the number of cells, repeated board letters, and
   several board paths spelling the same word.}
  {I build one Trie for all words so their shared prefixes are explored once.
   DFS advances through both the board and Trie, marks the current cell, and
   restores it during backtracking. Missing Trie edges provide the crucial
   pruning.}

\chaptercheckpoint{Search many words that share prefixes on one board.}
  {A Trie node synchronized with a marked board path.}
  {Running a completely separate unpruned DFS for every word.}{Trie build $O(W)$; search is prefix-pruned.}


\problemstart{Word Search}{79}{Medium}{Grid backtracking}{Matrix and recursion}

\subsection{The Problem}
Return whether one word can be formed from adjacent cells without reusing a
cell in the same path.

\begin{invariantbox}{the marked path spells the matched word prefix}
At recursion index \(i\), all earlier characters have been matched by distinct
adjacent cells. The current cell must equal character \(i\), after which it is
temporarily marked during all four choices.
\end{invariantbox}

\begin{visualbox}{Backtracking explores and restores a board path}
\centering
\begin{tikzpicture}[font=\small]
  \matrix[matrix of nodes,nodes={arraycell},row sep=-\pgflinewidth,
    column sep=-\pgflinewidth] (m) {
    |[fill=softorange,draw=orange,line width=1pt]| A &
    |[fill=softorange,draw=orange,line width=1pt]| B &
    |[fill=softorange,draw=orange,line width=1pt]| C & E\\
    S & F & |[fill=softorange,draw=orange,line width=1pt]| C & S\\
    A & |[fill=softorange,draw=orange,line width=1pt]| D &
    |[fill=softorange,draw=orange,line width=1pt]| E & E\\};
  \node[right=10mm of m-2-4,align=left]
    {matched path: \texttt{ABCCED}\\
     moves: right, right, down, down, left\\
     restore every cell when returning};
\end{tikzpicture}
\end{visualbox}

\subsection{Pseudocode}
\begin{lstlisting}[style=pseudocode]
for every board cell:
    if DFS(cell, word index 0): return true
DFS(cell, index):
    if index equals word length: return true
    reject invalid, visited, or mismatching cells
    mark cell; search four neighbors for index + 1
    restore cell and return whether any branch succeeded
return false
\end{lstlisting}

\subsection{C++ Implementation}
\begin{lstlisting}[style=compactcode,float=htbp,caption={Complete solution: Word Search}]
#include <string>
#include <vector>

using namespace std;

class Solution {
    bool dfs(vector<vector<char>>& b,const string& w,int i,int r,int c){
        if(i==(int)w.size()) return true;
        if(r<0||c<0||r==(int)b.size()||c==(int)b[0].size()||b[r][c]!=w[i]) return false;
        char ch=b[r][c]; b[r][c]='#';
        bool ok=dfs(b,w,i+1,r+1,c)||dfs(b,w,i+1,r-1,c)||dfs(b,w,i+1,r,c+1)||dfs(b,w,i+1,r,c-1);
        b[r][c]=ch; return ok;
    }
public:
    bool exist(vector<vector<char>>& board,string word){
        for(int r=0;r<(int)board.size();++r)for(int c=0;c<(int)board[0].size();++c)if(dfs(board,word,0,r,c))return true; return false;
    }
};
\end{lstlisting}

\subsection{Time and Space}
Worst-case time is \(O(mn\,4^L)\), with \(O(L)\) recursion space. Restoring
the character is required even when a branch fails.

\lessonexpansion
  {To find \texttt{ABCCED}, start at the matching \texttt{A}, mark it, and
   recursively follow adjacent \texttt{B}, \texttt{C}, \texttt{C},
   \texttt{E}, and \texttt{D}. A failed branch restores every marked cell so
   another route may use it.}
  {The DFS state records a board position and the next word index. It explores
   every legal non-reusing path with that prefix; marking prevents cycles and
   restoration keeps the board for sibling search branches.}
  {Check a one-letter word, a word longer than the board, repeated letters,
   paths requiring turns, a tempting path that reuses a cell, and no match.}
  {I start backtracking from every possible first character. Each recursive
   step matches one letter, temporarily marks the cell, explores four
   neighbors, and restores the cell afterward.}

\chaptercheckpoint{Match one word along a non-reusing grid path.}
  {A word index and temporarily marked board cells.}
  {Leaving failed-path cells marked.}{$O(mn\,4^L)$ worst-case time.}


\patternintoc{Binary Search and Subsequences}
\problemstart{Search in Rotated Sorted Array}{33}{Medium}{Modified binary search}{Array}

\subsection{The Problem}
Find a target in a rotated sorted array of distinct values and return its
index, or \(-1\).

\subsection{One Half Is Always Sorted}
At each midpoint, at least one side is normally ordered. Check whether the
target lies inside that sorted value range; keep that half if it does,
otherwise discard it.

\begin{invariantbox}{the target, if present, remains inside the boundaries}
The sorted-half range test identifies a half that cannot contain the target.
Discarding only that half keeps every possible target index.
\end{invariantbox}

\begin{visualbox}{Use the sorted half to choose safely}
\centering
\begin{tikzpicture}[font=\small]
  \foreach \x/\v in {0/4,1/5,2/6,3/7,4/0,5/1,6/2}{\node[arraycell] (a\x) at (1.0*\x,0) {\v};}
  \node[above=7mm of a3,text=orange]{mid};
  \node[draw=teal,very thick,fit=(a0)(a1)(a2)(a3),inner sep=1mm] {};
  \node[below=7mm of a1]{left half sorted};
\end{tikzpicture}
\end{visualbox}

\subsection{Pseudocode}
\begin{lstlisting}[style=pseudocode]
left = 0; right = n - 1
while left <= right:
    middle = midpoint
    if nums[middle] is target: return middle
    identify the sorted half
    if target lies inside that half: keep it
    else: keep the other half
return -1
\end{lstlisting}

\subsection{C++ Implementation}
\begin{lstlisting}[style=compactcode,float=htbp,caption={Complete solution: Search in Rotated Sorted Array}]
#include <vector>

using namespace std;

class Solution {
public:
    int search(vector<int>& nums,int target){
        int left=0,right=(int)nums.size()-1;
        while(left<=right){int mid=left+(right-left)/2;if(nums[mid]==target)return mid;
            if(nums[left]<=nums[mid]){
                if(nums[left]<=target&&target<nums[mid])right=mid-1;else left=mid+1;
            }else{
                if(nums[mid]<target&&target<=nums[right])left=mid+1;else right=mid-1;
            }
        } return -1;
    }
};
\end{lstlisting}

\subsection{Time and Space}
Each comparison removes half the range: \(O(\log n)\) time and \(O(1)\)
space. The distinct-value condition avoids ambiguous equal boundaries.

\lessonexpansion
  {For \texttt{[4,5,6,7,0,1,2]} and target 0, the first midpoint is value 7.
   The left half is sorted, but 0 is not in \([4,7)\), so discard it. In the
   remaining range, binary search reaches value 0 at index 4.}
  {At least one half around the midpoint is sorted. If the target lies inside
   that half's inclusive/exclusive value bounds, the opposite half cannot
   contain it; otherwise the sorted half can be discarded. So the target is
   keepd whenever it exists.}
  {Check no rotation, rotation by one position, target at either boundary,
   target at the pivot, two elements, one element, and a missing target.}
  {I first identify the sorted half, then use ordinary value-range reasoning
   to decide whether the target belongs there. This retains binary search's
   logarithmic elimination despite the rotation.}

\chaptercheckpoint{Binary search an array with one rotation break.}
  {Two boundaries and identification of the sorted half.}
  {Testing target indexes instead of sorted value ranges.}{$O(\log n)$ time.}


\problemstart{Longest Increasing Subsequence}{300}{Medium}{Subsequence DP}{Array}

\subsection{The Problem}
Return the length of the longest strictly increasing subsequence. Chosen
positions retain order but need not be adjacent.

\begin{invariantbox}{dp[i] is the best increasing subsequence ending at i}
Every previous \(j<i\) with \(nums[j]<nums[i]\) can extend its best
subsequence by the current value. Taking the maximum considers every possible
penultimate element.
\end{invariantbox}

\begin{visualbox}{A current value extends compatible earlier endings}
\centering
\begin{tikzpicture}[font=\small]
  \foreach \x/\v in {0/10,1/9,2/2,3/5,4/3,5/7,6/101,7/18}{
    \node[arraycell] (a\x) at (0.9*\x,0) {\v};
  }
  \foreach \x in {2,4,5,7}{
    \node[draw=teal,very thick,fit=(a\x),inner sep=-0.3pt] {};
  }
  \node[below=7mm of a4,text=teal,font=\bfseries]
    {read highlighted positions left to right: \(2<3<7<18\)};
\end{tikzpicture}
\end{visualbox}

\subsection{Pseudocode}
\begin{lstlisting}[style=pseudocode]
dp[i] = 1 for every position
for i from 0 to n-1:
    for j from 0 to i-1:
        if nums[j] < nums[i]:
            dp[i] = max(dp[i], dp[j] + 1)
return maximum dp value
\end{lstlisting}

\subsection{C++ Implementation}
\begin{lstlisting}[style=compactcode,float=htbp,caption={Complete solution: Longest Increasing Subsequence}]
#include <algorithm>
#include <vector>

using namespace std;

class Solution {
public:
    int lengthOfLIS(vector<int>& nums){
        vector<int> dp(nums.size(),1); int answer=1;
        for(int i=0;i<(int)nums.size();++i){
            for(int j=0;j<i;++j) if(nums[j]<nums[i]) dp[i]=max(dp[i],dp[j]+1);
            answer=max(answer,dp[i]);
        } return answer;
    }
};
\end{lstlisting}

\subsection{Time and Space}
The direct update rule uses \(O(n^2)\) time and \(O(n)\) space. A patience
sorting method improves time to \(O(n\log n)\), but this table most directly
shows the subsequence state.

\lessonexpansion
  {For \texttt{[10,9,2,5,3,7,101,18]}, every state begins at one. Value 5
   extends the subsequence ending at 2 to length two; value 7 extends either
   \(2,5\) or \(2,3\) to length three; value 18 extends a length-three state
   to the final answer four.}
  {Every increasing subsequence ending at index \(i\) has some penultimate
   index \(j<i\) with a smaller value, unless its length is one. Testing all
   such \(j\) includes the previous position in a best subsequence,
   making \(dp[i]\) correct.}
  {Check a decreasing array, an increasing array, all equal values, one value,
   repeated values mixed with growth, negative values, and a best subsequence
   that does not include the final element.}
  {I define \(dp[i]\) as the best increasing subsequence ending exactly at
   \(i\). Every smaller earlier value is a legal previous. The maximum over
   all ending positions is the answer.}

\chaptercheckpoint{Positions may be skipped while values strictly increase.}
  {The best subsequence length ending at each index.}
  {Using non-strict comparison and accepting equal values.}{$O(n^2)$ time and $O(n)$ space.}


\patternintoc{DP, Trees, and Reachability}
\problemstart{Climbing Stairs}{70}{Easy}{Fibonacci-style DP}{Rolling states}

\subsection{The Problem}
Count distinct ways to reach step \(n\) using moves of one or two steps.

\begin{invariantbox}{the two rolling values solve the previous two steps}
Every route to step \(i\) arrives from \(i-1\) or \(i-2\), and these cases are
separate, so their counts add.
\end{invariantbox}

\begin{visualbox}{Each step receives ways from the previous two steps}
\centering
\begin{tikzpicture}[font=\small,>={Stealth[length=2.2mm,width=1.55mm]}]
  \node[dsnode,visited] (s2) {step 2\\2 ways};
  \node[dsnode,active,below=10mm of s2] (s3) {step 3\\3 ways};
  \node[dsnode,right=35mm of $(s2)!0.5!(s3)$] (s4) {step 4\\5 ways};
  \draw[flowarrow,teal] (s2.east)--node[edgelabel,above]{take a 2-step move} (s4.west);
  \draw[flowarrow,orange] (s3.east)--node[edgelabel,below]{take a 1-step move} (s4.west);
  \node[right=7mm of s4,font=\bfseries] {\(2+3=5\)};
\end{tikzpicture}
\end{visualbox}

\subsection{Pseudocode}
\begin{lstlisting}[style=pseudocode]
if n <= 2: return n
oneBack = 2; twoBack = 1
for step from 3 to n:
    current = oneBack + twoBack
    twoBack = oneBack; oneBack = current
return oneBack
\end{lstlisting}

\subsection{C++ Implementation}
\begin{lstlisting}[style=compactcode,float=htbp,caption={Complete solution: Climbing Stairs}]
using namespace std;

class Solution {
public:
    int climbStairs(int n){
        int twoBack=1,oneBack=1;
        for(int step=2;step<=n;++step){int current=oneBack+twoBack;twoBack=oneBack;oneBack=current;}
        return oneBack;
    }
};
\end{lstlisting}

\subsection{Time and Space}
Time is \(O(n)\), space is \(O(1)\). The conceptual count for step zero is one
empty way, which makes the update rule work at the boundary.

\lessonexpansion
  {The counts for steps 0 through 5 are \(1,1,2,3,5,8\). For step 5, every
   valid climb ends with either a one-step move from step 4 or a two-step move
   from step 3, so the two preceding counts add to eight.}
  {The final move splits all valid climbs into two separate groups: those
   arriving from \(n-1\) and those arriving from \(n-2\). Their counts sum,
   and only the two previous states are required.}
  {Check \(n=1\), \(n=2\), the smallest allowed input, and larger values near
   the return-type limit.}
  {This is Fibonacci-style dynamic programming. I keep only the previous two
   counts because each new state depends exclusively on them.}

\chaptercheckpoint{Every state can be reached by one of two final moves.}
  {Counts for the previous two steps.}
  {Setting the zero-step base count to zero.}{$O(n)$ time and $O(1)$ space.}


\problemstart{Invert Binary Tree}{226}{Easy}{Recursive child swap}{Binary tree}

\subsection{The Problem}
Mirror a binary tree by exchanging the left and right child of every node.

\begin{invariantbox}{a completed call returns a fully mirrored subtree}
The current node swaps its child pointers, then recursively mirrors both
subtrees. Every edge below the node is reversed horizontally once.
\end{invariantbox}

\begin{visualbox}{Every left and right edge changes sides}
\centering
\begin{tikzpicture}[font=\small,>={Stealth[length=2.2mm,width=1.55mm]}]
  \node[trienode] (r1) {4}; \node[trienode,below left=10mm and 12mm of r1] (a1) {2};
  \node[trienode,below right=10mm and 12mm of r1] (b1) {7};
  \draw[-,thick,navy] (r1)--(a1) (r1)--(b1);
  \node[below=5mm of a1,xshift=12mm]{before};
  \node[trienode,right=52mm of r1] (r2) {4};
  \node[trienode,below left=10mm and 12mm of r2] (b2) {7};
  \node[trienode,below right=10mm and 12mm of r2] (a2) {2};
  \draw[-,thick,navy] (r2)--(b2) (r2)--(a2);
  \draw[flowarrow,orange] (r1.east)--node[edgelabel,above,fill=white]{swap children}(r2.west);
  \node[below=5mm of b2,xshift=12mm,text=teal,font=\bfseries]{after};
\end{tikzpicture}
\end{visualbox}

\subsection{Pseudocode}
\begin{lstlisting}[style=pseudocode]
invert(node):
    if node is null: return null
    swap node.left and node.right
    invert(node.left)
    invert(node.right)
    return node
\end{lstlisting}

\subsection{C++ Implementation}
\begin{lstlisting}[style=compactcode,float=htbp,caption={Complete solution: Invert Binary Tree}]
#include <utility>

using namespace std;

class Solution {
public:
    TreeNode* invertTree(TreeNode* root){
        if(root==nullptr) return nullptr;
        swap(root->left,root->right);
        invertTree(root->left); invertTree(root->right); return root;
    }
};
\end{lstlisting}

\subsection{Time and Space}
Every node is visited once: \(O(n)\) time and \(O(h)\) recursion space.

\lessonexpansion
  {For root 4 with children 2 and 7, swap them first. Recursively swap the
   children below 7 and below 2 as well. Applying the same operation a second
   time restores the original tree.}
  {A tree is mirrored exactly when every node exchanges its left and right
   subtrees and both exchanged subtrees are themselves mirrored. The recursive
   procedure performs this definition at every node.}
  {Check a null tree, one node, a skewed tree, an asymmetric tree, and verify
   that double inversion returns the original structure.}
  {I swap each node's child pointers and recursively invert both resulting
   subtrees. Every node is processed once, with stack depth equal to height.}

\chaptercheckpoint{Mirror every parent-child relationship.}
  {A subtree root whose children are swapped recursively.}
  {Swapping only the root's immediate children.}{$O(n)$ time and $O(h)$ space.}


\problemstart{Pacific Atlantic Water Flow}{417}{Medium}{Reverse multi-source DFS}{Grid and two reachability sets}

\subsection{The Problem}
Return cells from which water can flow to both oceans. Water moves to adjacent
cells of equal or lower height; the Pacific touches top and left edges, and
the Atlantic touches bottom and right edges.

\subsection{Search Backward from the Oceans}
Forward search from every cell repeats work. Reverse the edges: start from
each ocean boundary and move to neighbors of equal or greater height. Cells
reached by both reverse searches can flow to both oceans in the original
direction.

\begin{invariantbox}{reverse-reached cells can flow to the starting ocean}
Every reverse step climbs from a cell to a neighbor at least as high. Reversing
that step gives a legal downhill or level flow back to the ocean.
\end{invariantbox}

\begin{visualbox}{Two ocean floods climb inward and overlap}
\centering
\begin{tikzpicture}[font=\small]
  \matrix[matrix of nodes,nodes={arraycell,minimum width=9mm},
    row sep=-\pgflinewidth,column sep=-\pgflinewidth] (m) {
    |[fill=softteal,draw=teal]|1 & |[fill=softteal,draw=teal]|2 &
    |[fill=softteal,draw=teal]|2 & |[fill=softblue,draw=navy,line width=1.1pt,double]|3\\
    |[fill=softteal,draw=teal]|3 & |[fill=softteal,draw=teal]|2 &
    |[fill=softteal,draw=teal]|3 & |[fill=softblue,draw=navy,line width=1.1pt,double]|4\\
    |[fill=softteal,draw=teal]|2 & |[fill=softteal,draw=teal]|4 &
    |[fill=softblue,draw=navy,line width=1.1pt,double]|5 & |[fill=softorange,draw=orange]|3\\
    |[fill=softblue,draw=navy,line width=1.1pt,double]|6 &
    |[fill=softblue,draw=navy,line width=1.1pt,double]|7 &
    |[fill=softorange,draw=orange]|1 & |[fill=softorange,draw=orange]|4\\};
  \draw[very thick,teal] (m-1-1.north west)--(m-1-4.north east);
  \draw[very thick,teal] (m-1-1.north west)--(m-4-1.south west);
  \draw[very thick,orange] (m-4-1.south west)--(m-4-4.south east);
  \draw[very thick,orange] (m-1-4.north east)--(m-4-4.south east);
  \node[above=6mm of m-1-2,text=teal,font=\bfseries]{Pacific};
  \node[below=6mm of m-4-3,text=orange,font=\bfseries]{Atlantic};
  \node[right=11mm of m-2-4,align=left] {
    \textcolor{teal}{teal: Pacific only}\\
    \textcolor{orange}{orange: Atlantic only}\\
    \textcolor{navy}{double navy: both oceans}};
\end{tikzpicture}
\end{visualbox}

\subsection{Pseudocode}
\begin{lstlisting}[style=pseudocode]
run reverse DFS from all Pacific border cells
run reverse DFS from all Atlantic border cells
reverse DFS may move to an equal-or-higher neighbor
answer = every cell reached by both traversals
return answer
\end{lstlisting}

\subsection{C++ Implementation}
\begin{lstlisting}[style=compactcode,float=htbp,caption={Complete solution: Pacific Atlantic Water Flow}]
#include <vector>

using namespace std;

class Solution {
    void flow(const vector<vector<int>>& h,int r,int c,vector<vector<bool>>& seen){
        if(seen[r][c]) return; seen[r][c]=true;
        static const int d[5]={1,0,-1,0,1};
        for(int k=0;k<4;++k){int nr=r+d[k],nc=c+d[k+1];
            if(nr>=0&&nc>=0&&nr<(int)h.size()&&nc<(int)h[0].size()&&!seen[nr][nc]&&h[nr][nc]>=h[r][c]) flow(h,nr,nc,seen);
        }
    }
public:
    vector<vector<int>> pacificAtlantic(vector<vector<int>>& heights){
        int m=heights.size(),n=heights[0].size();
        vector<vector<bool>> pac(m,vector<bool>(n)),atl=pac;
        for(int r=0;r<m;++r){flow(heights,r,0,pac);flow(heights,r,n-1,atl);}
        for(int c=0;c<n;++c){flow(heights,0,c,pac);flow(heights,m-1,c,atl);}
        vector<vector<int>> out;
        for(int r=0;r<m;++r)for(int c=0;c<n;++c)if(pac[r][c]&&atl[r][c])out.push_back({r,c});
        return out;
    }
};
\end{lstlisting}

\subsection{Why It Works and Complexity}
Each ocean traversal visits each cell at most once, so time and reachability
space are \(O(mn)\); recursion can also reach \(O(mn)\). Intersecting the two
visited grids is needed: union would include cells reaching only one ocean.

\lessonexpansion
  {Seed the Pacific search from the top and left borders and the Atlantic
   search from the bottom and right. From an ocean cell, move only to neighbors
   of equal or greater height. A high interior cell reached by both searches
   appears in both Boolean grids and is added to the result.}
  {A reverse edge from height \(h\) to a neighbor of height at least \(h\)
   corresponds to a legal forward water step back toward the ocean. So each
   visited grid contains exactly the cells that can reach its ocean. Their
   intersection is exactly the requested answer.}
  {Check one cell, one row, one column, all equal heights, strictly increasing
   terrain, strictly decreasing terrain, plateaus, and cells lying directly on
   both ocean boundaries.}
  {Instead of starting a repeated downhill search from every cell, I reverse
   the graph and launch two multi-source traversals from the ocean borders.
   Cells visited by both traversals can flow to both oceans.}

\chaptercheckpoint{Many cells ask whether downhill paths reach two boundaries.}
  {Two reverse reachability grids seeded from ocean edges.}
  {Moving downhill during the reverse search or taking the union.}{$O(mn)$ time and space.}
